% !TEX root = ../main.tex
\section{Introduction}
\label{sec:intro}

Wigner's 1961 thought experiment~\cite{wigner1961} exposes a structural tension at the core of quantum theory: an observer (``the Friend'') sealed inside an isolated laboratory registers a definite measurement outcome, whereas an external superobserver (Wigner) modeling the laboratory unitarily assigns an entangled superposition to the joint system-plus-friend state. Modern extensions combining Wigner's setup with bipartite Bell correlations---initiated by Frauchiger and Renner~\cite{frauchiger2018} and Brukner~\cite{brukner2018}, tested in photonic proof-of-principle architectures~\cite{proietti2019,bong2020}, and sharpened into the device-independent \emph{Local Friendliness} (LF) theorem by Bong \emph{et al.}~\cite{bong2020} and subsequent generalizations~\cite{wiseman2023,haddara2023}---prove that quantum predictions violate inequalities derived solely from \emph{Local Agency} and the \emph{Absoluteness of Observed Events} (AOE), without assuming predetermination or outcome independence. Parallel advances on superconducting quantum processors have demonstrated contextuality-based quantum--classical separations in tasks spanning up to 105 qubits~\cite{kumar2026} and begun probing multi-qubit EWFS circuits directly~\cite{zeng2024}, highlighting both the viability of coherent quantum registers as controlled test beds and the rapid degradation of coherent advantage under circuit depth and environmental coupling.

Despite their foundational force, existing no-go theorems leave two concrete gaps between abstract correlation polytopes and relativistic physics:
\begin{enumerate}
    \item \textbf{Covariant Localization of Relative Facts:} Rejecting AOE avoids nonlocal collapse, yet relational approaches risk ambiguity unless the boundary separating determinate facts from indeterminate propositions is anchored covariantly to spacetime causal structure and equipped with a consistency law when previously separated observers compare records.
    \item \textbf{Quantitative Leakage Scaling Across Genuine LF Facets:} Idealized EWFS derivations assume unit coherence retention ($f = 1$) during Wigner's reversal unitary $U^\dagger$. In realistic processors, partial distinguishability $\ell = \sqrt{1 - f^2} > 0$ inevitably leaks into unreversed modes. For a single leaky laboratory the optimum is known: it is the trade-off between CHSH violation and an eavesdropper's guessing probability established for device-independent randomness certification~\cite{pironio2010,masanes2011,acin2012}, with the environment in the role of the eavesdropper, and for two qubits it also follows from the Horodecki criterion~\cite{horodecki1995}. Partial classical records~\cite{baumann2023} and imperfect record objectivity~\cite{rivlin2025} are likewise known to interpolate between violation and non-violation of LF inequalities. What has been missing is an account of how leakage in both laboratories affects LF polytope facets that are not facets of the Bell polytope.
\end{enumerate}

In this work, we address both gaps through a unified algebraic and operational framework. In Sec.~\ref{sec:horizons}, we formulate the \emph{Causal Actuality Hypothesis} (CAH), which assigns truth valuations $v_O(P, \tau) \in \{T, F, U\}$ relative to stabilized record subalgebras $\mathcal{R}_O(\tau) \subset \mathfrak{A}(J^-(x_O(\tau)))$ inside an observer's causal past light cone, building on the relativistic algebraic measurement program of Fewster and Verch~\cite{fewster2020} and Alves~\cite{alves2024}. While three-valued quantum valuations trace back to Reichenbach~\cite{reichenbach1944}, CAH deploys $\{T, F, U\}$ strictly as covariant causal bookkeeping on local algebras rather than altering classical logic on shared records. This perspective refines Relational Quantum Mechanics~\cite{rovelli1996,smerlak2007,dibiagio2021} by replacing bare pairwise interaction events with causal record horizons. It complements Quantum Darwinism analyses of Wigner's friend~\cite{zurek2018,rivlin2025}, which already constrain LF violations through the non-objectivity of the friend's record~\cite{rivlin2025}, with closed-form bounds for a pure-dephasing leakage model.

In Sec.~\ref{sec:polytope_bounds}, we combine exact polytope verification in $\mathbb{R}^{15}$ with operator-algebraic bounds. First, we give a Sum-of-Squares certificate, valid in arbitrary unital $C^*$-algebras, that the one-friend Semi-Brukner inequality obeys $B_{\max}^{\mathrm{SB}}(f) = 2\sqrt{1 + f^2} = 2\sqrt{2 - \ell^2}$; this restates the randomness trade-off of Refs.~\cite{pironio2010,masanes2011,acin2012} for the EWFS dephasing model. Second, we analyze bilateral friend dephasing in two-qubit scenarios across the five canonical $(3,3,2,2)$ inequality categories tested by Bong \emph{et al.}~\cite{bong2020}, uncovering a striking separation between quadratic one-friend violation excess $\mathcal{O}(f^2)$ and quartic two-friend suppression $\mathcal{O}(f^4)$ as $f \to 0$; for two-qubit CHSH-type functionals we prove that the excess is governed by the product of the two laboratories' coherence retentions through the bound $2\sqrt{1 + f_A^2 f_B^2}$, which for the Brukner functional holds in every dimension, whereas for the other functionals the quartic law is specific to complete friend measurements on qubits. In Sec.~\ref{sec:dsw_ceilings}, we address John Bell's critique of ``For All Practical Purposes'' (FAPP) solutions~\cite{bell1990}: recognizing that flat-spacetime quadrupole gravitational radiation can be suppressed via adiabatic uncomputation, we anchor the fundamental limit on reversible control to unshieldable soft-graviton and soft-photon flux across Killing and cosmological de Sitter horizons established by Danielson, Satishchandran, and Wald~\cite{danielson2023,danielson2025}, and discuss composite sub-horizons and open experimental questions.
